\documentclass[border=10pt]{standalone}
\usepackage[utf8]{inputenc}
\usepackage[T1]{fontenc}
\usepackage{tabularx}
\usepackage{booktabs}
\usepackage{enumitem}
\usepackage{array}
\usepackage{lmodern}

\newcommand{\usecasetable}[9]{
    \renewcommand{\arraystretch}{1.4}
    \begin{tabularx}{19cm}{|l|>{\raggedright\arraybackslash}X|}
        \hline
        \textbf{ID Use Case} & #1 \\ \hline
        \textbf{Nome Use Case} & #2 \\ \hline
        \textbf{Attori} & #3 \\ \hline
        \textbf{Descrizione} & #4 \\ \hline
        \textbf{Precondizioni} & #5 \\ \hline
        \textbf{Postcondizioni} & #6 \\ \hline
        \textbf{Flusso Principale} & #7 \\ \hline
        \textbf{Flussi alternativi} & #8 \\ \hline
        \textbf{Business Rules e Riferimenti} & #9 \\ \hline
    \end{tabularx}
}

\begin{document}
	\usecasetable
	{UC-07}
	{Manage Reviews}
	{Customer}
	{Permette ad un Customer di valutare e commentare la propria esperienza presso una spiaggia in cui ha effettivamente soggiornato, oppure di eliminare una propria recensione precedentemente pubblicata.}
	{L'utente è loggato. La spiaggia è in stato ATTIVO. L'utente non ha un Ban per quella spiaggia. Deve esistere almeno una prenotazione per quella spiaggia con data nel passato (rispetto a quando viene applicato lo Use Case) e in stato CONFIRMED.}
	{La recensione viene salvata o rimossa dal database, associata alla spiaggia e all'utente.}
	{
		\begin{enumerate}[nosep, leftmargin=*]
			\item L'utente naviga nella sezione "Le mie prenotazioni" o sulla pagina della spiaggia;
			\item L'utente clicca sul pulsante "Lascia una recensione" per la spiaggia selezionata;
			\item Il sistema verifica che tutti i requisiti siano soddisfatti (prenotazione passata e CONFIRMED, assenza di ban, spiaggia attiva);
			\item Il sistema mostra il modulo della recensione;
			\item L'utente seleziona una valutazione (da 1 a 5 stelle) e inserisce un commento testuale;
			\item L'utente clicca su "Pubblica";
			\item Il sistema salva la recensione nel database e mostra un messaggio di conferma.
		\end{enumerate}
	}
	{
		\begin{enumerate}[nosep, leftmargin=25px]
			\item[\textbf{2a.}] \textbf{Eliminazione Recensione:} 
			\begin{enumerate}[nosep, leftmargin=*]
				\item[i.] L'utente seleziona una recensione che ha già pubblicato in passato e clicca su "Elimina";
				\item[ii.] Il sistema chiede conferma dell'operazione ("Sei sicuro di voler eliminare questa recensione?");
				\item[iii.] L'utente conferma;
				\item[iv.] Il sistema rimuove il record dal database e mostra un messaggio di successo.
			\end{enumerate}
			\item[\textbf{3a-e.}] \textbf{Mancanza Requisiti:} Se l'utente tenta di accedere all'URL della recensione senza averne diritto (es. prenotazione futura, prenotazione CANCELLED/REJECTED o mai stato in quella spiaggia), il sistema blocca l'azione mostrando: "Devi aver completato un soggiorno in questa struttura per poter lasciare una recensione.".
			\item[\textbf{3b-e.}] \textbf{Spiaggia Disattivata/Ban:} Se la spiaggia non è più attiva o l'utente è stato bannato da quella struttura dopo il soggiorno, il sistema inibisce la funzione di recensione.
			\item[\textbf{6a-e.}] \textbf{Recensione non completata:} Se l'utente non compila entrambi i campi, il sistema evidenzia i campi da compilare obbligatoriamente.
		\end{enumerate}
	}
	{
		\begin{enumerate}[nosep, leftmargin=*]
			\item[•] La valutazione è obbligatoria ed è un valore intero tra 1 e 5.
			\item[•] Il commento testuale è anch'esso obbligatorio.
			\item[•] Il sistema accetta recensioni solo per prenotazioni la cui data è strettamente minore della data odierna (es. prenotazione il 19 luglio, recensibile dal 20 luglio).
			\item[•] L'App Ban è già gestito a monte dal Login (UC-01).
		\end{enumerate}
		
		\begin{enumerate}[leftmargin=*]
			\item[•] \textbf{Mockup}: si veda Figura 17 - Le Mie Prenotazioni e Recensioni
			
			\item[•] \textbf{Testing}: le classi che si occupano dei test sono \texttt{ReviewServiceTest} nel package \texttt{service}; \texttt{ReviewTest} nel package \texttt{domain}.
		\end{enumerate}
	}
\end{document}